\section{Sovereign Access Control}

\label{sec:access}

\subsection{Three-Tier Governance Model}

ZDC implements three access governance tiers reflecting the diversity of data sensitivity in biodiversity research:

\begin{table}[H]
\centering
\caption{Access governance tiers.}
\label{tab:access_tiers}
\begin{tabularx}{\textwidth}{@{}llX@{}}
\toprule
\textbf{Tier} & \textbf{Default Access} & \textbf{Use Cases} \\
\midrule
Open & Unrestricted & Published occurrence records, aggregated statistics, metadata \\
Permissioned & Request-based & Sensitive location data (endangered species), pre-publication data \\
Sovereign & Community-governed & Traditional ecological knowledge, indigenous land observations \\
\bottomrule
\end{tabularx}
\end{table}

\subsection{Permissioned Access}

Permissioned data uses attribute-based encryption (ABE) where the decryption key is conditioned on verifiable attributes of the requester:

\begin{equation}
    \text{Enc}(\mathcal{D}, \text{policy}) \rightarrow \text{CT}
\end{equation}

\begin{equation}
    \text{Dec}(\text{CT}, \text{key}_{\text{attributes}}) \rightarrow \begin{cases} \mathcal{D} & \text{if attributes satisfy policy} \\ \perp & \text{otherwise} \end{cases}
\end{equation}

Policies are expressed as Boolean formulas over attributes, such as:

\begin{verbatim}
    (role:researcher AND affiliation:university)
    OR (role:government AND jurisdiction:national)
    OR (approved-by:data-owner)
\end{verbatim}

\subsection{Sovereign Data Governance}

The sovereign tier implements the CARE Principles for Indigenous Data Governance \citep{carroll2020care}.
Each sovereign dataset is governed by a \emph{Data Governance Council} (DGC), a multisignature authority comprising community-designated representatives:

\begin{equation}
    \text{Access}(\text{ZDO}_{\text{sovereign}}) \leftarrow \text{Vote}_{\text{DGC}}(\text{request}) \geq \text{threshold}
\end{equation}

The DGC can:
\begin{itemize}[leftmargin=*]
    \item Set access policies with cultural context annotations (e.g., seasonal restrictions on sacred site data).
    \item Require benefit-sharing agreements as a condition of access.
    \item Attach Traditional Knowledge (TK) labels specifying culturally appropriate use.
    \item Revoke access retroactively if terms are violated.
    \item Delegate decision-making to subcommittees for specific data categories.
\end{itemize}

\subsection{Spatial Obfuscation}

For sensitive species (e.g., those vulnerable to poaching), ZDC supports spatial obfuscation that degrades location precision while preserving analytical utility:

\begin{equation}
    (x', y') = \text{Obfuscate}(x, y, r) = (x + \eta_x, y + \eta_y) \quad \text{where } \|\boldsymbol{\eta}\| \leq r
\end{equation}

The obfuscation radius $r$ is species-specific, determined by threat level (IUCN Red List status) and local poaching risk assessment.
We implement a Geo-Indistinguishability mechanism based on the Laplace distribution:

\begin{equation}
    \Pr[(x', y') \mid (x, y)] \propto e^{-\epsilon \cdot d((x', y'), (x, y))}
\end{equation}

\noindent where $\epsilon$ controls the privacy-utility tradeoff.
Researchers needing precise locations can request permissioned access through the standard mechanism.

% ============================================================
% 6. Curation Market
% ============================================================
